\metatitle{Kromatic symmetric functions}
\metadescription{An introduction to kromatic symmetric functions, their definition as a K-theoretical analog of chromatic symmetric functions, and their properties.}


\section[kromaticSymmetric]{Kromatic symmetric functions}


\begin{polydata}{kromaticSymmetric}
  Name   & Kromatic symmetric functions \\
  Space  & Sym \\
  Basis  & False \\
  Rating & 2 \\
  Bib    & CrewPechenikSpirkl2023x \\
  Year   & 2023 \\
  KTheoreticAnalogueOf & chromaticQuasisymmetric | CrewPechenikSpirkl2023x \\
  PositiveIn & grothendieckStable | CrewPechenikSpirkl2023x | scope=claw-free incomparability graphs of posets \\
\end{polydata}


A $K$-theoretical analog of the chromatic symmetric functions were introduced in \cite{CrewPechenikSpirkl2023x},
called the \defin{kromatic symmetric functions}.
These serve as a non-homogeneous lift of the \hyperref[chromaticQuasisymmetricDefinition]{chromatic symmetric functions} introduced by \name{Richard Stanley} \cite{Stanley1995}.
One main feature is that these expand positively into \hyperref[grothendieckStable]{Grothendieck symmetric functions} 
whenever the graph is a claw-free incomparability graph of a poset.

\name{Eric Marberg} gives a Hopf-algebraic construction of the kromatic
symmetric function and proves that it has a positive expansion into
multifundamental quasisymmetric functions \cite[Cor.~3.25]{Marberg2025}.
He also introduces two quasisymmetric $q$-analogues.  One is
multifundamental-positive but is symmetric only for cluster graphs
\cite[Thms.~4.7 and~4.8]{Marberg2025}; the other has a positive expansion
into stable Grothendieck functions when the graph is the incomparability
graph of a natural unit interval order \cite[Thm.~4.21]{Marberg2025}.

\name{Laura Pierson} gives an explicit expansion of the kromatic symmetric
function in the $K$-analogue of the power-sum basis and proves integrality
and a sign rule for the coefficients \cite[Thm.~1.1 and
Cor.~1.5]{Pierson2026}.  A later refinement gives a Lyndon-heap
interpretation of these coefficients and proves that the
$\omega$-image is positive in the same $K$-power-sum basis
\cite[Thms.~1.2 and~1.3]{Pierson2025}.

The kromatic symmetric function contains more information than the ordinary
chromatic symmetric function.  In particular, knowing it is equivalent to
knowing the multiset of independence polynomials of all induced subgraphs
\cite[Cor.~1.6]{Pierson2025}, and it determines the number of induced copies
of several small independence-unique graphs \cite{Pierson2024x}.  However, it
does not distinguish all graphs: \name{Laura Pierson} and \name{Soham Samanta}
construct four pairs of non-isomorphic graphs on eight vertices with the same
kromatic symmetric function \cite{PiersonSamanta2026}.
